% Contenu partagé (élève/corrigé) — IParcours G2
% !TeX program = lualatex

% Barème indicatif (ajustez si besoin)
\bareme{
  \exbadge{1}{5} \quad
  \exbadge{2}{4} \quad
  \exbadge{3}{3}
}

% =============================
% EXERCICE 1 — Cours et propriétés (8 pts)
% =============================
\begin{exercisebox}{\faIcon{book} Exercice 1 : Cours et propriétés \points{5}}
  \textbf{1.} Énonce le \textbf{théorème de Pythagore}. \points{1}

  \anslines[0.7cm]{3}{\begin{minipage}{\linewidth}
    Dans un triangle rectangle, le carré de la longueur de l'hypoténuse est égal à la somme des carrés des longueurs des deux autres côtés.
  \end{minipage}}

  \smallspace

  \textbf{2.} Complète chaque propriété avec l’égalité de Pythagore correspondante. \points{2}
  \begin{enumerate}[label=\alph*)]
    \item Si un triangle $FQZ$ est rectangle en $Q$ alors\;\onlyeleve{\underline{\hspace{2.2cm}} + \underline{\hspace{2.2cm}} = \underline{\hspace{2.2cm}}} \onlycorrige{$\;FZ^2 = FQ^2 + QZ^2$}.
    \item Si un triangle $JMX$ est rectangle en $X$ alors\;\onlyeleve{\underline{\hspace{2.2cm}} + \underline{\hspace{2.2cm}} = \underline{\hspace{2.2cm}}} \onlycorrige{$\;JM^2 = JX^2 + XM^2$}.
    \item Si un triangle $KNS$ est rectangle en $K$ alors\;\onlyeleve{\underline{\hspace{2.2cm}} + \underline{\hspace{2.2cm}} = \underline{\hspace{2.2cm}}} \onlycorrige{$\;NS^2 = NK^2 + KS^2$}.
    \item Si un triangle $PIO$ est rectangle en $I$ alors\;\onlyeleve{\underline{\hspace{2.2cm}} + \underline{\hspace{2.2cm}} = \underline{\hspace{2.2cm}}} \onlycorrige{$\;PO^2 = PI^2 + IO^2$}.
  \end{enumerate}

  \smallspace

  \textbf{3.} Énonce la \textbf{réciproque du théorème de Pythagore}. \points{1}

  \anslines[0.7cm]{3}{\begin{minipage}{\linewidth}
    Si, dans un triangle, le carré de la longueur du plus grand côté est égal à la somme des carrés des longueurs des deux autres côtés, alors ce triangle est rectangle (l’angle droit est à l’extrémité commune aux deux petits côtés).
  \end{minipage}}

  \smallspace

  \textbf{4.} Complète chaque propriété (utiliser la réciproque). \points{1}
  \begin{enumerate}[label=\alph*)]
    \item Si $ST^2 + TJ^2 = SJ^2$ alors le triangle \onlyeleve{\underline{\hspace{1.6cm}}} \onlycorrige{$\;STJ$} est rectangle en \onlyeleve{\underline{\hspace{1cm}}} \onlycorrige{$\;T$}.
    \item Si $MR^2 + MV^2 = RV^2$ alors le triangle \onlyeleve{\underline{\hspace{1.6cm}}} \onlycorrige{$\;MRV$} est rectangle en \onlyeleve{\underline{\hspace{1cm}}} \onlycorrige{$\;M$}.
    \item Si $WX^2 + XY^2 = WY^2$ alors le triangle \onlyeleve{\underline{\hspace{1.6cm}}} \onlycorrige{$\;WXY$} est rectangle en \onlyeleve{\underline{\hspace{1cm}}} \onlycorrige{$\;X$}.
  \end{enumerate}
\end{exercisebox}

\exspace

% =============================
% EXERCICE 2 — Vérifier si un triangle est rectangle (4 pts)
% =============================
\begin{exercisebox}{\faIcon{check-double} Exercice 2 : Triangle rectangle ? \points{4}}
  \textbf{7.} Soit $BUT$ un triangle tel que $TU = \SI{5,1}{cm}$, $BU = \SI{7,9}{cm}$ et $BT = \SI{6}{cm}$. Ce triangle est-il rectangle ? Justifie ta réponse. \points{2}

  \anslines[0.55cm]{6}{\begin{minipage}{\linewidth}
    $\max(BU, BT, TU) = BU$. On compare $BU^2$ et $BT^2 + TU^2$ :\\
    $BU^2 = 7{,}9^2 = 62{,}41$,\quad $BT^2 + TU^2 = 6^2 + 5{,}1^2 = 36 + 26{,}01 = 62{,}01$.\\
    Comme $62{,}41 \neq 62{,}01$, on n’a pas l’égalité de Pythagore. Donc le triangle $BUT$ n’est \textbf{pas} rectangle.
  \end{minipage}}

  \smallspace

  \textbf{8.} Soit $CAR$ un triangle tel que $RA = \SI{3,3}{cm}$, $AC = \SI{5,6}{cm}$ et $RC = \SI{6,5}{cm}$. Ce triangle est-il rectangle ? Justifie ta réponse. \points{2}

  \anslines[0.55cm]{6}{\begin{minipage}{\linewidth}
    $\max(RC, RA, AC) = RC$. On compare $RC^2$ et $RA^2 + AC^2$ :\\
    $RA^2 + AC^2 = 3{,}3^2 + 5{,}6^2 = 10{,}89 + 31{,}36 = 42{,}25$;\quad $RC^2 = 6{,}5^2 = 42{,}25$.\\
    Les deux quantités sont égales, donc par la réciproque de Pythagore, $CAR$ est rectangle en $A$.
  \end{minipage}}
\end{exercisebox}

\exspace

% =============================
% EXERCICE 3 — Problème (vie courante) : Table de chevet (3 pts)
% =============================
\begin{problembox}{\faIcon{cubes} Exercice 3 : La table de chevet \points{3}}
  Un designer a imaginé une table de chevet en marbre (plateau) et métal (structure). On modélise la structure ci-dessous. Les segments \([JL]\) et \([KI]\) sont \textbf{perpendiculaires} en \(M\).

  \begin{center}
    \begin{tikzpicture}[scale=0.12]
      % Coins du cadre
      \coordinate (J) at (0,40);
      \coordinate (I) at (64,40);
      \coordinate (L) at (64,0);
      \coordinate (K) at (0,0);
      % Intersection des renforts (choisi comme (12.8,12.8) transformé ?
      % On place M pour respecter les demi-longueurs sur les renforts :
      % JK et IL ne sont volontairement pas utilisées pour fixer M ;
      % on règle M pour avoir JM=38.4 sur JL et IM=51.2 sur KI, etc.
      % Configuration géométrique : M = (28.8, 1*?) mais on annote sans exiger l'échelle.
      % On trace seulement une figure illustrative (non à l'échelle exacte).

      % Cadre
      \draw[very thick,secondary!80] (J) -- (I) -- (L) -- (K) -- cycle;
      % Renforts
      \path (J) -- (L) coordinate[pos=0.64] (M1); % pour annoter 38,4 et 21,6
      \path (K) -- (I) coordinate[pos=0.36] (M2); % pour annoter 28,8 et 51,2
      % On force un seul point M visuel (proche de l’intersection)
      \coordinate (M) at ($(M1)!0.5!(M2)$);
      \draw[thick] (J) -- (L);
      \draw[thick] (K) -- (I);
      % Annotation longueurs le long des renforts
      \draw (J) -- (M) node[midway,sloped,above] {\small \SI{38,4}{cm}};
      \draw (M) -- (L) node[midway,sloped,below] {\small \SI{21,6}{cm}};
      \draw (K) -- (M) node[midway,sloped,above] {\small \SI{28,8}{cm}};
      \draw (M) -- (I) node[midway,sloped,above] {\small \SI{51,2}{cm}};
      % Angles droits au centre
      \draw ($(M)+(-2,0)$) -- ($(M)+(-2,2)$) -- ($(M)+(0,2)$);
      % Légendes
      \node[above left]  at (J) {\small J};
      \node[above right] at (I) {\small I};
      \node[below right] at (L) {\small L};
      \node[below left]  at (K) {\small K};
      \node[below]       at (M) {\small M};
      % Plateau
      \node[above] at ($(J)!0.5!(I)$) {\footnotesize Plateau en marbre};
    \end{tikzpicture}
  \end{center}

  \textbf{Questions.}
  \begin{enumerate}[label=\arabic*.]
    \item Calcule les longueurs \(JK\) et \(JI\). \points{2}
      \anslines[0.6cm]{4}{\onlycorrige{\(JK = \sqrt{JM^2+MK^2} = \sqrt{38{,}4^2+28{,}8^2} = \sqrt{2304} = \SI{48}{cm}.\\
      JI = \sqrt{JM^2+MI^2} = \sqrt{38{,}4^2+51{,}2^2} = \sqrt{4096} = \SI{64}{cm}.\)}}

    \item En déduire si le montant \([JK]\) est vertical et si le plateau \([JI]\) est parallèle au sol. \points{1}
      \anslines[0.6cm]{4}{\onlycorrige{On a \(KI = KM + MI = 28{,}8 + 51{,}2 = \SI{80}{cm}.\)
      Dans le triangle \(JKI\), \(JK^2 + JI^2 = 48^2 + 64^2 = 2304 + 4096 = 6400 = 80^2 = KI^2\). \\
      Donc l’angle \(\widehat{KJI}\) est droit: \(JK \perp JI\). Ainsi, \(JK\) est vertical et \(JI\) est horizontal (donc parallèle au sol).}}
  \end{enumerate}
\end{problembox}
